% ===============================================================
% FILE: main_tesis.tex
% DOCUMENTO DE TESIS MAGÍSTER (FORMATO INSTITUCIONAL USM)
% Objetivo: Ser exhaustivo, cubriendo cada detalle metodológico y su justificación biológica.
% ===============================================================

\documentclass[a4paper, 12pt]{book}
% --- Paquetes de LaTeX Esenciales ---
\usepackage[utf8]{inputenc}
\usepackage[T1]{fontenc}
\usepackage{amsmath, amssymb} % Para ecuaciones matemáticas
\usepackage{graphicx} % Para incluir figuras
\usepackage{caption} 
\usepackage{booktabs} % Para tablas profesionales (toprule, midrule...)
\usepackage{amsthm} % Para definiciones y teoremas

% --- Configuración de estilo e idioma ---
\usepackage[spanish]{babel}
\setcounter{secnumdepth}{2} % Numerar hasta la subsección

% ... [Otros paquetes requeridos: bibliography, amsmath, etc.] ... 


\title{Desarrollo y Validación de Técnicas Avanzadas de Grafos para el Procesamiento EEG.}
\author{Tu Nombre Completo\\Mentor Académico (USM)}
\date{\today}

\begin{document}

\maketitle

% --- SECCIONES PRELIMINARES ---
\chapter*{Agradecimientos} % [Aquí va texto en español]
\addcontentsline{toc}{chapter}{Agradecimientos}

\chapter*{Resumen Abstracto (Español)} % Debe ser la sinopsis más concisa.
\addcontentsline{toc}{chapter}{Resumen Abstracto (Español)}

% -----------------------------------

\tableofcontents % LaTeX genera el índice automáticamente

\chapter{Marco Teórico de las Señales Biofisiológicas}
\label{cap:marco_teorico}
% Contenido de la Revisión Bibliográfica. Aquí se usa Citas [2], [3], [5].
% Se debe crear una subsección crítica sobre por qué ICA/Spline es insuficiente.

\chapter{Graph Signal Processing (GSP): Fundamentos Matemáticos}
\label{cap:gsp_teorico}
% Aquí se define el grafo W, Laplaciano L=D-W y los operadores de filtrado en el dominio del grafo. 
% Se deben incluir ecuaciones matemáticas detalladas para la construcción de W (Kernel vs Data-Driven).

\chapter{Metodología Experimental: El Sistema EARS}
\label{cap:metodologia_experimental}
Este capítulo describe nuestra arquitectura reproducible, dividida en tres fases que garantizan la validez estadística.

\section{Fase I: Construcción del Grafo $\mathbf{W}$}
\label{sec:grafos_w}
Detalla el proceso comparativo de $\text{GRA-04}$ vs $\text{GRA-08}$, incluyendo las ecuaciones y los parámetros clave ($\theta, \kappa, \sigma$).

\section{Fase II: El Protocolo EARS (La Máquina del Tiempo)}
\label{sec:pipeline_arquitectura}
Aquí se debe describir el diagrama de flujo completo. Se narra cómo la señal pasa por las múltiples tareas (filtrado $\rightarrow$ imputación $\rightarrow$ rechazo).

\subsection{2.1 Protocolo de Datos Faltantes ($\text{PRT-01}$)} 
Detallar el proceso `Orchestrator` y el uso del $ground\text{-}truth$.

\section{Fase III: Evaluación Estadística Rigurosa}
(Aquí va la explicación detallada sobre los $\text{QA Gates}$: Manifiesta que no solo se corre, sino que se *valida* si es correcto.)

% ===============================================================
% Contenido para: Cap. IV Discusión Científica 
% Objetivo: Interpretar los datos y justificar la superioridad de GSP-TVT.
% ===============================================================

\section{Discusión de Resultados}
\label{sec:discussion_results}
El conjunto de resultados acumulados en las diversas iteraciones (ver \ref{fig:tv_vs_instant_family}) valida nuestra hipótesis central: la modelación del espacio tiempo-conectividad ($\mathbf{W}_{sct}$) es fundamental para el preprocesamiento EEG.

\subsection{Superioridad Estadística y la Limitación del Modelo Lineal}
La mejora estadísticamente significativa de MAE (media $\pm$ desviación) en la familia TV/Tiempo, como se observó notablemente en las iteraciones \texttt{it124} et al., indica que los métodos lineales son incapaces de modelar el fenómeno biológico del tiempo. La magnitud de esta mejora ($\text{p}=0.027$, $\text{Gain}>+40\%$) supera la mera corrección de datos faltantes; es una reestructuración fundamental en el entendimiento de las dependencias temporales corticales.

\subsection{El Rol Crítico de la Trazabilidad y Robustez}
Un logro metodológico clave fue la creación del $\text{Orchestrator}$ (Protocolo EARS). Al validar los resultados con escenarios controlados \emph{y} aleatorios ($\text{PRT-01}$) en múltiples datasets (MNE, BCI), se demostró que el rendimiento es robusto a las condiciones de adquisición. La capacidad de cuantificar y mitigar la varianza mediante el sistema $\text{QA Gate}$ asegura que los hallazgos son generalizables y no un artefacto del *dataset* específico.

\subsection{Conclusiones Condicionales: El Horizonte Clínico}
El resultado final nos obliga a una **Declaración de Dominio de Aplicabilidad**. La superioridad $\text{TV/Tiempo}$ es indiscutible en el contexto de la dinámica cerebral, pero nuestro análisis indica que métodos como $\text{nnk}$ o $\text{spherical\_spline}$ son más adecuados cuando la varianza temporal no es un factor primario (ej., detección simple de artefactos). Por lo tanto, nuestra contribución principal no es un "mejor método", sino una **matriz de selección de método óptimo** basada en el contexto clínico/investigativo.

% ===========================================
% Contenido para: Cap. V - Conclusión (La última palabra)
% Objetivo: Resumir el impacto sin repetir los resultados; abrir puertas a futuro.
% ===========================================

\section{Conclusiones}
\label{sec:conclusion_final}
Este trabajo culmina la propuesta de un marco avanzado y reproducible para el preprocesamiento EEG, superando las limitaciones inherentes a los modelos clásicos mediante la integración del Graph Signal Processing en el dominio espacio-tiempo. Las contribuciones fundamentales incluyen:

\begin{itemize}
    \item \textbf{Diseño Sistemático:} La creación del sistema EARS, un protocolo de validación exhaustiva que asegura la trazabilidad completa (Protocolo $\text{B3/B4}$).
    \item \textbf{Modelado Avanzado:} Demostramos que el modelo $\text{TV/Tiempo}$ es superior porque eleva el modelado a una dinámica espacio-temporal.
    \item \textbf{Guía de Aplicabilidad:} Ofrecemos una matriz de decisión basada en el dominio del problema, lo cual permite seleccionar la mejor herramienta según si el fenómeno primario es espacial ($\text{ICA}$) o temporal ($\text{GSP}_{\text{TV}}$).
\end{itemize}

\section{Trabajo Futuro (El Motor para los Próximos 5 Años)}
(Aquí muestras que piensas más allá de la tesis.)
Las futuras investigaciones deben centrarse en:
\begin{enumerate}
    \item \textbf{Causalidad:} Pasar de la correlación ($\mathbf{W}$) a modelos causales. Esto implica integrar análisis como el $\text{Granger Causality}$ sobre las señales reconstruidas por GSP.
    \item \textbf{Multimodalidad:} Integrar otras fuentes de datos (Imágenes, Movimiento) para crear grafos híbridos y aún más robustos.
    \item \textbf{Aprendizaje Adaptativo:} Desarrollar modelos $\text{ML}$ que aprendan dinámicamente los hiperparámetros óptimos ($\theta$, $\kappa$) en tiempo real durante la adquisición de datos, eliminando la intervención manual del investigador.
\end{enumerate}
